\PuzzleChapterIntro{The Tangent-Wrought Five-Count Seal}{Titan \textbullet\ Valuations \textbullet\ Theorem Web \textbullet\ Geometry Thread}{This is a test of structural arithmetic. The numbers involved grow far too large for direct computation; you must navigate the web of theorems that governs their prime factorizations and valuations.}
\Lore{
A chalked circle is scratched into the brass: \textit{radius by squares; direction by tangent}.

Behind it sits a five-toothed wheel, heavy in the palm---each tooth stamped with a recurrence instead of a numeral.

A second inscription runs along the axle, shallow under lamplight:
\textit{When the reader alternates, only the last tooth speaks.}

The seal accepts one number and refuses the rest.
}

\Problem
\begingroup\footnotesize

For a nonzero integer \(x\), let \(\nu(x)\) be the largest \(t\ge 0\) such that \(2^t\mid x\), and set \(\nu(x)=\nu(|x|)\).

\smallskip
\noindent\textbf{Geometry key.}
In the coordinate plane, let
\[
\mathcal C:\ x^2+y^2=25,\qquad \mathcal L:\ x+y=7.
\]
Let \(P\) be the intersection point of \(\mathcal C\) and \(\mathcal L\) in the first quadrant with \(y>x\).
Let \(\ell\) be the tangent line to \(\mathcal C\) at \(P\), and write its slope in lowest terms as
\[
\operatorname{slope}(\ell)=-\frac{u}{v},\qquad u,v\in\mathbb Z_{>0},\ \gcd(u,v)=1.
\]
Define the \emph{scale} \(\kappa=v^2\) and the \emph{count} \(m=OP\) (distance from the origin to \(P\); in this configuration \(m\) is an integer).

\smallskip
\noindent\textbf{The teeth.}
Define a sequence \((t_j)_{j\ge 0}\) by
\[
t_0=u,\qquad t_{j+1}=t_j^2-2t_j+2\quad (j\ge 0),
\]
and form the set of \(m\) scaled teeth
\[
T=\{\kappa t_0,\kappa t_1,\dots,\kappa t_{m-1}\}.
\]

\smallskip
\noindent\textbf{The alternating reader.}
For any subset \(A\subseteq T\), list its elements in \emph{strictly decreasing} order \(a_1>\cdots>a_p\) and define
\[
r(A)=\sum_{k=1}^{p}(-1)^{k-1}a_k,\qquad r(\varnothing)=0.
\]
Let
\[
R=\sum_{A\subseteq T} r(A),\qquad e=\nu(R),\qquad n=e+1.
\]

\smallskip
\noindent\textbf{Two auxiliary locks.}
Define \((b_k)_{k\ge 0}\) by
\[
b_0=0,\quad b_1=1,\quad b_k=2b_{k-1}+b_{k-2}\ (k\ge 2),
\]
and set
\[
B=2^{\nu(b_{2^n})},\qquad A=2^{\nu(u^{2^n}-1)}.
\]
Also define
\[
D=\gcd\!\left(u^{2^n}-1,\ 2^{2^m}-1\right).
\]

\smallskip
\noindent\textbf{The clasp.}
Let \(W\) be the (unique, if it exists) power of \(2\) that can be written as
\[
W=(a+b^2)(b+a^2)\quad\text{for some }a,b\in\mathbb Z_{>0}.
\]

\smallskip
\noindent\textbf{Seal-number.}
\[
\mathcal N \;=\; 2^{e}\cdot \frac{A}{B}\cdot W\cdot D.
\]

\smallskip
\endgroup

\VaultDirective{Determine the exact value of $\mathcal N$.}

\SolutionPage
\scriptsize
\textbf{1) Geometry: \((u,v,\kappa,m)\).}
From \(x^2+(7-x)^2=25\) we get \(x\in\{3,4\}\); with \(y>x\) take \(P=(3,4)\).
The radius \(OP\) has slope \(4/3\), so the tangent has slope \(-3/4\). Hence
\[
u=3,\ v=4,\ \kappa=v^2=16,\ \text{and }m=OP=\sqrt{3^2+4^2}=5\in\mathbb Z.
\]


\textbf{2) Teeth: Fermat ladder and product.}
From \(t_{j+1}=t_j^2-2t_j+2\),
\[
t_0=3,\ t_1=5,\ t_2=17,\ t_3=257,\ t_4=65537.
\]
Also \(t_{j+1}-2=t_j(t_j-2)\), so by induction
\[
t_0t_1\cdots t_{k-1}=t_k-2\qquad(k\ge 1).
\]
Since \(t_{j+1}=(t_j-1)^2+1\) and \(t_0=2^{2^0}+1\), we have \(t_j=2^{2^j}+1\), so \(t_5=2^{32}+1\) and
\[
t_0t_1t_2t_3t_4=t_5-2=2^{32}-1.
\]
The same identity gives \(\gcd(t_i,t_j)=1\) for \(i\ne j\).


\textbf{3) Alternating reader: only the maximum survives.}
For \(X=\{x_1<\cdots<x_m\}\),
$
\sum_{A\subseteq X} r(A)=2^{m-1}x_m.
$
Indeed, fixing \(x_i\), its total contribution is \(x_i\cdot 2^{i-1}(1-1)^{m-i}\), since there are \(2^{i-1}\) choices of smaller elements and the sign depends only on chosen larger elements. This vanishes unless \(i=m\).

Here \(T\) has \(m=5\) elements and \(\max(T)=\kappa t_4=16\cdot 65537\), so
\[
R=2^{4}\cdot 16\cdot 65537=2^{8}\cdot 65537,\quad e=\nu(R)=8,\quad n=e+1=9.
\]


\textbf{4) Hinge valuation: \(\nu(b_k)=\nu(k)\).}
Write \((1+\sqrt2)^k=U_k+b_k\sqrt2\) with integers \(U_k,b_k\). Then \(b_k\) satisfies the given recurrence and \(U_k\) is odd.
Squaring gives \(b_{2k}=2U_kb_k\), hence \(\nu(b_{2k})=1+\nu(b_k)\). Since \(b_k\) is odd iff \(k\) is odd, repeated halving yields \(\nu(b_k)=\nu(k)\). Therefore
\[
\nu(b_{2^n})=n,\qquad B=2^{\nu(b_{2^n})}=2^n=2^9=512.
\]


\textbf{5) Dyadic depth of \(3^{2^n}-1\).}
For \(n\ge 1\) we claim \(3^{2^n}-1\equiv 2^{n+2}\pmod{2^{n+3}}\), hence \(\nu(3^{2^n}-1)=n+2\). For \(n=1\), \(3^2-1=8\equiv 2^3\pmod{16}\). If \(3^{2^n}\equiv 1+2^{n+2}\pmod{2^{n+3}}\), then squaring gives \(3^{2^{n+1}}\equiv 1+2^{n+3}\pmod{2^{n+4}}\).
With \(n=9\),
\[
\nu(3^{512}-1)=11,\qquad A=2^{11}=2048,\qquad \frac{A}{B}=\frac{2048}{512}=4.
\]

\textbf{6) The clasp.}
If \((a+b^2)(b+a^2)=2^t\), then \(a+b^2=2^m\) and \(b+a^2=2^k\).
Since \(2^m\) is even, \(a\equiv b\pmod 2\), so \(a+b-1\) is odd. WLOG \(m\ge k\); subtracting gives
\[
2^m-2^k=(b-a)(a+b-1).
\]
Because \(a+b-1\) is odd, \(\nu(2^m-2^k)=k\) forces \(2^k\mid (b-a)\).
But \(0\le b-a<b<2^k\) (since \(b+a^2=2^k>b\)), so \(b=a\).
Then \(a+b^2=a+a^2=a(a+1)\) is a power of \(2\), so \(a\) and \(a+1\) are both powers of \(2\), hence \(a=1\).
Thus \(W=4\).

\textbf{7) The latch \(D=\gcd(3^{512}-1,2^{32}-1)\).}
By Step 2, \(2^{32}-1=t_0t_1t_2t_3t_4\) with coprime factors, so
\[
D=\prod_{0\le i\le 4,\ t_i\mid(3^{512}-1)} t_i.
\]
Now:
\[
3\nmid(3^{512}-1),\qquad 3^4\equiv 1\pmod 5, \ 3^8\equiv -1\pmod{17}\ \Rightarrow\ 3^{16}\equiv 1\pmod{17},
\] 
\[
3^8\equiv 136,\ 3^{16}\equiv -8,\ 3^{32}\equiv 64,\ 3^{64}\equiv -16,\ 3^{128}\equiv -1\pmod{257}\ \Rightarrow\ 3^{256}\equiv 1\pmod{257},
\]
and repeated squaring modulo \(65537\) yields \(3^{16}\equiv -11088\), \(3^{32}\equiv 61869\), \(3^{64}\equiv 19139\), \(3^{128}\equiv 15028\), \(3^{256}\equiv 282\pmod{65537}\). Since \(282\not\equiv \pm 1\pmod{65537}\), we have \(3^{512}\not\equiv 1\pmod{65537}\), so \(65537\nmid(3^{512}-1)\).
Hence
\[
D=5\cdot 17\cdot 257=21845.
\]

\textbf{8) Assemble.}
\[
\mathcal N=2^e\cdot \frac{A}{B}\cdot W\cdot D
=256\cdot 4\cdot 4\cdot 21845
=4096\cdot 21845
=89{,}477{,}120.
\]

\FinalAnswer{89477120}


\NextPage
